\documentclass[../../main.tex]{subfiles}

\begin{document}

\section{实数的开集、闭集和 Borel 集}\label{section:ch01-开集闭集Borel}

\begin{definition}[开集]\label{definition:ch01-开集}
一个实数的集合 $\mathcal{O}$ 称为开的, 若对每个 $x\in\mathcal{O}$, 存在 $r>0$ 使得区间 $(x-r,x+r)$ 包含于 $\mathcal{O}$.
\end{definition}

对于 $a<b$, 区间 $(a,b)$ 是一个开集. 事实上, 令 $x$ 属于 $(a,b)$. 定义 $r=\min\{b-x,x-a\}$. 观察到 $(x-r,x+r)$ 包含于 $(a,b)$. 因此 $(a,b)$ 是开有界区间, 且每个有界开区间都是这种形式. 对于 $a,b\in\mathbf{R}$, 我们定义
\begin{align*}
(a,\infty)=\{x\in\mathbf{R}\mid a<x\},\quad (-\infty,b)=\{x\in\mathbf{R}\mid x<b\},\quad (-\infty,\infty)=\mathbf{R}.
\end{align*}
观察到每个这样的集合是一个开区间. 此外, 不难看出, 由于每个实数集在扩充实数集中有下确界与上确界, 每个无界开区间都是上述形式.

\begin{proposition}\label{proposition:ch01-开集性质}
实数集 $\mathbf{R}$ 和空集 $\varnothing$ 是开的, 任何开集的有限族的交是开的, 任何开集族的并是开的.
\end{proposition}

\begin{proof}
显然 $\mathbf{R}$ 和 $\varnothing$ 是开的, 而任何开集族的并是开的. 令 $\{\mathcal{O}_{k}\}_{k=1}^{n}$ 为 $\mathbf{R}$ 的开子集的有限族. 若该族的交是空的, 则交是空集, 因此是开的. 否则, 令 $x$ 属于 $\bigcap_{k=1}^{n}\mathcal{O}_{k}$. 对于 $1\leqslant k\leqslant n$, 选取 $r_{k}>0$ 使得 $(x-r_{k},x+r_{k})\subseteq\mathcal{O}_{k}$. 定义 $r=\min\{r_{1},\cdots,r_{n}\}$. 则 $r>0$ 且 $(x-r,x+r)\subseteq\bigcap_{k=1}^{n}\mathcal{O}_{k}$. 因此 $\bigcap_{k=1}^{n}\mathcal{O}_{k}$ 是开的.
\end{proof}

然而, 任何开集族的交是开的不成立. 例如, 对每个自然数 $n$, 令 $\mathcal{O}_{n}$ 为开区间 $(-1/n,1/n)$. 则根据 $\mathbf{R}$ 的 Archimedeas 性质, $\bigcap_{n=1}^{\infty}\mathcal{O}_{n}=\{0\}$, 而 $\{0\}$ 不是一个开集.

\begin{proposition}\label{proposition:ch01-开区间分解}
每个非空开集是可数个不交开区间族的并.
\end{proposition}

\begin{proof}
令 $\mathcal{O}$ 为 $\mathbf{R}$ 的非空开子集. 令 $x$ 属于 $\mathcal{O}$. 存在 $y>x$ 使得 $(x,y)\subseteq\mathcal{O}$, 且存在 $z<x$ 使得 $(z,x)\subseteq\mathcal{O}$. 定义扩充的实数 $a_{x}$ 和 $b_{x}$ 为
\begin{align*}
a_{x}=\inf\{z\mid (z,x)\subseteq\mathcal{O}\}\quad\text{与}\quad b_{x}=\sup\{y\mid (x,y)\subseteq\mathcal{O}\}.
\end{align*}
则 $I_{x}=(a_{x},b_{x})$ 是包含 $x$ 的开区间. 我们宣称
\begin{align}
I_{x}\subseteq\mathcal{O}\text{ 但 }a_{x}\notin\mathcal{O},\quad b_{x}\notin\mathcal{O}.\label{align:ch01-区间分解}
\end{align}
事实上, 令 $w$ 属于 $I_{x}$, 比如 $x<w<b_{x}$. 根据 $b_{x}$ 的定义, 存在数 $y>w$ 使得 $(x,y)\subseteq\mathcal{O}$, 因而 $w\in\mathcal{O}$. 此外, $b_{x}\notin\mathcal{O}$, 因为若 $b_{x}\in\mathcal{O}$, 则对某个 $r>0$ 我们有 $(b_{x}-r,b_{x}+r)\subseteq\mathcal{O}$. 因此 $(x,b_{x}+r)\subseteq\mathcal{O}$, 与 $b_{x}$ 的定义矛盾. 类似地, $a_{x}\notin\mathcal{O}$. 考虑开区间族 $\{I_{x}\}_{x\in\mathcal{O}}$. 由于 $\mathcal{O}$ 中的每个 $x$ 是 $I_{x}$ 的成员, 而每个 $I_{x}$ 包含于 $\mathcal{O}$, 我们有 $\mathcal{O}=\bigcup_{x\in\mathcal{O}}I_{x}$. 我们从\eqref{align:ch01-区间分解}推出 $\{I_{x}\}_{x\in\mathcal{O}}$ 是不交的. 因此 $\mathcal{O}$ 是不交的开区间族的并. 剩下来要证明该族是可数的. 根据有理数的稠密性（\cref{theorem:ch01-有理数稠密}）, 这些开区间的每一个包含一个有理数. 这建立了开区间族与有理数子集之间的一一对应. 我们从\cref{theorem:ch01-可数子集}和\cref{corollary:ch01-可数推论}推出任何有理数集是可数的. 因此 $\mathcal{O}$ 是可数个不交开区间族的并.
\end{proof}

\begin{definition}[闭包点、闭包、闭集]\label{definition:ch01-闭包闭集}
对于实数集 $E$, $x$ 称为 $E$ 的闭包点, 若每个包含 $x$ 的开区间也包含 $E$ 的点. $E$ 的全体闭包点称为 $E$ 的闭包且记为 $\overline{E}$.

显然我们总是有 $E\subseteq\overline{E}$. 若 $E$ 包含它的所有闭包点, 即 $E=\overline{E}$, 则集合 $E$ 称为闭的.
\end{definition}

\begin{proposition}\label{proposition:ch01-闭包性质}
对于实数集 $E$, 它的闭包 $\overline{E}$ 是闭的. 此外, $\overline{E}$ 在以下意义下是包含 $E$ 的最小闭集: 若 $F$ 是闭的且 $E\subseteq F$, 则 $\overline{E}\subseteq F$.
\end{proposition}

\begin{proof}
集合 $\overline{E}$ 是闭的, 若它包含所有闭包点. 令 $x$ 为 $\overline{E}$ 的闭包点. 考虑包含 $x$ 的开区间 $I_{x}$. 存在一个点 $x'\in\overline{E}\cap I_{x}$. 由于 $x'$ 是 $E$ 的闭包点, 且开区间 $I_{x}$ 包含 $x'$, 存在点 $x''\in E\cap I_{x}$. 因此每个包含 $x$ 的开区间也包含 $E$ 的点, 且因此 $x\in\overline{E}$. 所以集合 $\overline{E}$ 是闭的. 显然, 若 $A\subseteq B$, 则 $\overline{A}\subseteq\overline{B}$, 因此, 若 $F$ 是闭的且包含 $E$, 则 $\overline{E}\subseteq\overline{F}=F$.
\end{proof}

\begin{proposition}\label{proposition:ch01-开闭互补}
实数集是开的当且仅当它在 $\mathbf{R}$ 中的补是闭的.
\end{proposition}

\begin{proof}
首先假定 $E$ 是 $\mathbf{R}$ 的一个开子集. 令 $x$ 为 $\mathbf{R}\setminus E$ 的闭包点. 则 $x$ 不属于 $E$, 因为否则就会有一个包含 $x$ 且包含于 $E$ 的开区间, 因而与 $\mathbf{R}\setminus E$ 不交. 于是 $x$ 属于 $\mathbf{R}\setminus E$ 且因此 $\mathbf{R}\setminus E$ 是闭的. 现在假定 $\mathbf{R}\setminus E$ 是闭的. 令 $x$ 属于 $E$. 则必有包含 $x$ 且包含于 $E$ 的开区间, 否则每个包含 $x$ 的开区间包含 $\mathbf{R}\setminus E$ 的点, 且因此 $x$ 是 $\mathbf{R}\setminus E$ 的闭包点. 由于 $\mathbf{R}\setminus E$ 是闭的, $x$ 也属于 $\mathbf{R}\setminus E$. 这是一个矛盾.
\end{proof}

由于 $\mathbf{R}\setminus[\mathbf{R}\setminus E]=E$, 从前一个命题得出一个集合是闭的当且仅当它的补是开的. 因此, 根据 De Morgan 等式, \cref{proposition:ch01-开集性质} 可用闭集重述如下.

\begin{proposition}\label{proposition:ch01-闭集性质}
空集 $\varnothing$ 和 $\mathbf{R}$ 是闭的, 任何闭集的有限族的并是闭的, 任何闭集族的交是闭的.
\end{proposition}

集族 $\{E_{\lambda}\}_{\lambda\in\Lambda}$ 称为是集合 $E$ 的覆盖, 若 $E\subseteq\bigcup_{\lambda\in\Lambda}E_{\lambda}$. 谈到 $E$ 的覆盖的子覆盖, 我们指的是该覆盖的子族自身也是 $E$ 的一个覆盖. 若覆盖中的每个集合 $E_{\lambda}$ 是开的, 我们称 $\{E_{\lambda}\}_{\lambda\in\Lambda}$ 为 $E$ 的一个开覆盖. 若覆盖 $\{E_{\lambda}\}_{\lambda\in\Lambda}$ 仅包含有限个集合, 我们称它为有限覆盖. 该术语是不一致的: "开覆盖" 中的 "开" 指的是该覆盖的集合; "有限覆盖" 中的 "有限" 指的是族而不是隐含该族中的集合是有限集. 因此, 术语 "开覆盖" 是语言的误用, 而恰当的说法应该是 "用开集覆盖". 遗憾的是, 前一个术语已在数学中广泛使用.

\begin{theorem}[Heine-Borel 定理]\label{theorem:ch01-HeineBorel}
令 $F$ 为闭有界实数集. 则 $F$ 的每个开覆盖有一个有限子覆盖.
\end{theorem}

\begin{proof}
我们首先考虑 $F$ 是闭有界区间 $[a,b]$ 的情形. 令 $\mathcal{F}$ 为 $[a,b]$ 的开覆盖. 定义 $E$ 为具有如下性质的区间 $[a,x]$, 即可被 $\mathcal{F}$ 的有限个集合覆盖的数 $x\in[a,b]$ 的集合. 由于 $a\in E$, $E$ 是非空的. 由于 $E$ 有上界 $b$, 根据 $\mathbf{R}$ 的完备性, $E$ 有上确界. 定义 $c=\sup E$. 由于 $c$ 属于 $[a,b]$, 存在 $\mathcal{O}\in\mathcal{F}$ 包含 $c$. 由于 $\mathcal{O}$ 是开的, 存在 $\varepsilon>0$, 使得区间 $(c-\varepsilon,c+\varepsilon)$ 包含于 $\mathcal{O}$. 现在 $c-\varepsilon$ 不是 $E$ 的上界, 因而必有 $x\in E$ 满足 $x>c-\varepsilon$. 由于 $x\in E$, 存在覆盖 $[a,x]$ 的 $\mathcal{F}$ 中的集合的有限族 $\{\mathcal{O}_{1},\cdots,\mathcal{O}_{k}\}$. 因此, 有限族 $\{\mathcal{O}_{1},\cdots,\mathcal{O}_{k},\mathcal{O}\}$ 覆盖区间 $[a,c+\varepsilon)$. 于是 $c=b$, 否则 $c<b$ 且 $c$ 不是 $E$ 的上界. 因此 $[a,b]$ 可被 $\mathcal{F}$ 中的有限个集合覆盖, 这证明了我们考虑的特殊情形.

现在令 $F$ 为任何闭有界集, 而 $\mathcal{F}$ 是 $F$ 的一个开覆盖. 由于 $F$ 是有界的, 它包含于某个有界闭区间 $[a,b]$. 前一个命题告诉我们集合 $\mathcal{O}=\mathbf{R}\setminus F$ 是开的, 因为 $F$ 是闭的. 令 $\mathcal{F}'$ 为添加 $\mathcal{O}$ 到 $\mathcal{F}$ 后得到的开集族, 即 $\mathcal{F}'=\mathcal{F}\cup\mathcal{O}$. 由于 $\mathcal{F}$ 覆盖 $F$, $\mathcal{F}'$ 覆盖 $[a,b]$. 根据我们刚考虑的情形, 存在 $\mathcal{F}'$ 的有限子族覆盖 $[a,b]$, 因此也覆盖 $F$. 通过从 $F$ 的这个有限子覆盖去掉 $\mathcal{O}$, 若 $\mathcal{O}$ 属于该有限子覆盖, 我们得到 $\mathcal{F}$ 中覆盖 $F$ 的有限族.
\end{proof}

我们说集合的可数族 $\{E_{n}\}_{n=1}^{\infty}$ 是下降的, 若对每个自然数 $n$, $E_{n+1}\subseteq E_{n}$. 说它是上升的, 若对每个自然数 $n$, $E_{n}\subseteq E_{n+1}$.

\begin{theorem}[集套定理]\label{theorem:ch01-集套}
令 $\{F_{n}\}_{n=1}^{\infty}$ 为下降的非空闭实数集的可数族, 其中 $F_{1}$ 有界. 则
\begin{align*}
\bigcap_{n=1}^{\infty}F_{n}\neq\varnothing.
\end{align*}
\end{theorem}

\begin{proof}
我们用反证法. 假定交集是空的. 则对每个实数 $x$, 存在自然数 $n$ 使得 $x\notin F_{n}$, 即 $x\in\mathcal{O}_{n}=\mathbf{R}\setminus F_{n}$. 因此 $\bigcup_{n=1}^{\infty}\mathcal{O}_{n}=\mathbf{R}$. 根据\cref{proposition:ch01-开闭互补}, 由于每个 $F_{n}$ 是闭的, 每个 $\mathcal{O}_{n}$ 是开的. 因此 $\{\mathcal{O}_{n}\}_{n=1}^{\infty}$ 是 $\mathbf{R}$ 的一个开覆盖, 从而也是 $F_{1}$ 的开覆盖. Heine-Borel 定理告诉我们存在自然数 $N$ 使得 $F_{1}\subseteq\bigcup_{n=1}^{N}\mathcal{O}_{n}$. 由于 $\{F_{n}\}_{n=1}^{\infty}$ 是下降的, 补集族 $\{\mathcal{O}_{n}\}_{n=1}^{\infty}$ 是上升的. 因此 $\bigcup_{n=1}^{N}\mathcal{O}_{n}=\mathcal{O}_{N}=\mathbf{R}\setminus F_{N}$. 因此 $F_{1}\subseteq\mathbf{R}\setminus F_{N}$, 这与 $F_{N}$ 是 $F_{1}$ 的非空子集的假设矛盾.
\end{proof}

\begin{definition}[$\sigma$ 代数]\label{definition:ch01-sigma代数}
给定集合 $X$, $X$ 的子集族 $\mathcal{A}$ 称为（$X$ 的子集的）$\sigma$ 代数, 若: (i) 空集 $\varnothing$ 属于 $\mathcal{A}$; (ii) $\mathcal{A}$ 中的集合在 $X$ 中的补也属于 $\mathcal{A}$; (iii) $\mathcal{A}$ 中集合的可数族的并也属于 $\mathcal{A}$.
\end{definition}

给定集合 $X$, 族 $\{\varnothing,X\}$ 是一个 $\sigma$ 代数, 它有两个成员且它包含于每个 $X$ 的子集的 $\sigma$ 代数. 另一个极端情形是 $X$ 的所有子集组成的集族且包含每个 $X$ 的子集的 $\sigma$ 代数 $2^{X}$. 对任何 $\sigma$ 代数 $\mathcal{A}$, 我们从 De Morgan 等式推出 $\mathcal{A}$ 关于属于 $\mathcal{A}$ 的集合的可数族的交封闭. 此外, 由于空集属于 $\mathcal{A}$, $\mathcal{A}$ 关于属于 $\mathcal{A}$ 的集合的有限并与有限交封闭. 我们也观察到 $\sigma$ 代数关于相对补封闭, 若 $A_{1}$ 和 $A_{2}$ 属于 $\mathcal{A}$, $A_{1}\setminus A_{2}=A_{1}\cap[X\setminus A_{2}]$ 也如此. 以下命题的证明直接从 $\sigma$ 代数的定义得到.

\begin{proposition}\label{proposition:ch01-最小sigma代数}
令 $\mathcal{F}$ 为集合 $X$ 的子集族. 则所有包含 $\mathcal{F}$ 的 $X$ 的子集的 $\sigma$ 代数的交 $\mathcal{A}$ 是一个包含 $\mathcal{F}$ 的 $\sigma$ 代数. 此外, 在任何包含 $\mathcal{F}$ 的 $\sigma$ 代数也包含 $\mathcal{A}$ 的意义下, 它是包含 $\mathcal{F}$ 的最小的 $X$ 的子集的 $\sigma$ 代数.
\end{proposition}

令 $\{A_{n}\}_{n=1}^{\infty}$ 为属于 $\sigma$ 代数 $\mathcal{A}$ 的集合的可数族. 由于 $\mathcal{A}$ 关于可数交与并封闭, 以下两个集合属于 $\mathcal{A}$:
\begin{align*}
\lim\sup\{A_{n}\}_{n=1}^{\infty}=\bigcap_{k=1}^{\infty}\left[\bigcup_{n=k}^{\infty}A_{n}\right]\quad\text{与}\quad\lim\inf\{A_{n}\}_{n=1}^{\infty}=\bigcup_{k=1}^{\infty}\left[\bigcap_{n=k}^{\infty}A_{n}\right].
\end{align*}
集合 $\lim\sup\{A_{n}\}_{n=1}^{\infty}$ 是对可数无穷多个指标 $n$ 属于 $A_{n}$ 的点的集合, 而集合 $\lim\inf\{A_{n}\}_{n=1}^{\infty}$ 是除指定至多有限多个指标 $n$ 外属于 $A_{n}$ 的点的集合.

虽然任何开集族的并是开的而有限个开集的交是开的, 如同我们看到的, 但是可数个开集的交不一定是开的. 在建立实直线上的 Lebesgue 测度与积分论时, 我们将看到包含开集的实数集的最小 $\sigma$ 代数是自然的研究对象.

\begin{definition}[Borel 集族]\label{definition:ch01-Borel集}
实数的 Borel 集族 $\mathcal{B}$ 是包含所有实数的开集的实数集的最小 $\sigma$ 代数.
\end{definition}

每个开集是 Borel 集, 且由于 $\sigma$ 代数关于补是封闭的, 我们从\cref{proposition:ch01-开闭互补}推出每个闭集是 Borel 集. 因此, 由于每个单点集是闭的, 每个可数集是 Borel 集. 开集的可数交称为 $G_{\delta}$ 集. 闭集的可数并称为 $F_{\sigma}$ 集. 由于 $\sigma$ 代数关于可数并与可数交封闭, 每个 $G_{\delta}$ 集和每个 $F_{\sigma}$ 集是 Borel 集. 此外, 开的或者闭的实数集的可数族的 $\lim\inf$ 和 $\lim\sup$ 都是 Borel 集.

\begin{exercise}\label{exercise:ch01-27}
有理数集是开集还是闭集?
\end{exercise}

\begin{exercise}\label{exercise:ch01-28}
哪些实数集既开又闭?
\end{exercise}

\begin{exercise}\label{exercise:ch01-29}
找到两个集合 $A$ 和 $B$ 使得 $A\cap B=\varnothing$ 而 $\overline{A}\cap\overline{B}\neq\varnothing$.
\end{exercise}

\begin{exercise}\label{exercise:ch01-30}
点 $x$ 称为集合 $E$ 的聚点, 若它是 $E\setminus\{x\}$ 的闭包点.
\begin{enumerate}[(i)]
\item 证明 $E$ 的聚点的集合 $E'$ 是一个闭集.
\item 证明 $\overline{E}=E\cup E'$.
\end{enumerate}
\end{exercise}

\begin{exercise}\label{exercise:ch01-31}
点 $x$ 称为集合 $E$ 的孤立点, 若存在 $r>0$ 使得 $(x-r,x+r)\cap E=\{x\}$. 证明: 若集合由孤立点组成, 则它是可数的.
\end{exercise}

\begin{exercise}\label{exercise:ch01-32}
点 $x$ 称为集合 $E$ 的内点, 若存在 $r>0$ 使得开区间 $(x-r,x+r)$ 包含于 $E$. $E$ 的内点的集合称为 $E$ 的内部, 记为 $\operatorname{int}E$. 证明:
\begin{enumerate}[(i)]
\item $E$ 是开的当且仅当 $E=\operatorname{int}E$.
\item $E$ 是稠密的当且仅当 $\operatorname{int}(\mathbf{R}\setminus E)=\varnothing$.
\end{enumerate}
\end{exercise}

\begin{exercise}\label{exercise:ch01-33}
证明: 若 $F_{1}$ 是无界的, 集套定理是错的.
\end{exercise}

\begin{exercise}\label{exercise:ch01-34}
证明 Heine-Borel 定理等价于实数的完备性公理. 证明集套定理等价于实数的完备性公理.
\end{exercise}

\begin{exercise}\label{exercise:ch01-35}
证明 Borel 集族是包含闭集的最小 $\sigma$ 代数.
\end{exercise}

\begin{exercise}\label{exercise:ch01-36}
证明 Borel 集族是包含形如 $[a,b)$（其中 $a<b$）的区间的最小 $\sigma$ 代数.
\end{exercise}

\begin{exercise}\label{exercise:ch01-37}
证明每个开集是一个 $F_{\sigma}$ 集.
\end{exercise}

\end{document}